% 04-compilation.tex

\section{Compilation and Execution}\label{sec:compilation}

Memory planning and Tier~1 replay are implemented.
SMT kernel synthesis and shadow handle execution are designs.

\subsection{Static Memory Planning}

Crucible replaces per-operation dynamic allocation with a static plan computed from the dataflow graph.
During \code{build\_trace()}, the background thread identifies unique tensor storages as \code{TensorSlot} structs (40 bytes each: \code{offset\_bytes}, \code{nbytes}, \code{birth\_op}, \code{death\_op}, dtype, device, \code{is\_external} flag, \code{slot\_id}).
Storage size is computed via \code{compute\_storage\_nbytes()}, which handles negative strides (e.g., \code{torch.flip}, \code{as\_strided}) by tracking both max and min offset contributions.
External tensors (parameters, data loader outputs) are excluded from the pool.

A sweep-line algorithm assigns offsets within a single contiguous pool using counting sort $O(n{+}k)$ rather than insertion sort:

\begin{enumerate}
  \item Sort tensors by birth index.
  \item Maintain free intervals within the pool.
  \item For each tensor in birth order: find the first fitting interval, assign offset, mark occupied until death.
  \item At death points, return intervals to the free list.
\end{enumerate}

The output is a \code{MemoryPlan}: \code{pool\_bytes} (total), \code{num\_slots}, \code{num\_external}, device context (\code{DeviceType}, \code{device\_idx}, \code{device\_capability}), and distributed topology (rank, world\_size).
At runtime, one allocation obtains the pool.
Every subsequent ``allocation'' is \code{base\_ptr + offset}.
No mutex, no freelist, no fragmentation.

\textbf{Correctness.}
Non-overlap is proved in Lean~4 for the general case and verified by Z3 for configurations up to 64 tensors with arbitrary lifetimes.
Figure~\ref{fig:memplan} illustrates the sweep-line assignment.

\begin{figure}[ht]
\centering
\begin{tikzpicture}[xscale=0.72, yscale=0.85]
  %% Grid
  \fill[black!2] (0,0) rectangle (12,4.2);
  \draw[black!10, very thin] foreach \x in {1,...,11} {(\x,0) -- (\x,4.2)};
  \draw[black!10, very thin] foreach \y in {1,...,4}  {(0,\y) -- (12,\y)};

  %% Axes
  \draw[black!40, -{Stealth[length=4pt]}] (0,0) -- (12.5,0)
    node[right, font=\sffamily\scriptsize] {op index};
  \draw[black!40, -{Stealth[length=4pt]}] (0,0) -- (0,4.6)
    node[above, font=\sffamily\scriptsize] {pool offset (KB)};

  %% Op ticks
  \foreach \x in {0,...,11} {
    \node[font=\tiny, text=black!40, below=1pt] at (\x+0.5, 0) {\x};
  }
  %% Offset ticks
  \foreach \y/\lab in {0/0, 1/1, 2/2, 3/3, 4/4} {
    \node[font=\tiny, text=black!40, left=2pt] at (0, \y) {\lab};
  }

  %% Tensor bars
  %% T0: ops 0-5, offset 0, size 2KB
  \fill[blue!18] (0.1, 0.05) rectangle (5.9, 1.95);
  \draw[blue!50, thick] (0.1, 0.05) rectangle (5.9, 1.95);
  \node[font=\sffamily\tiny\bfseries, blue!70] at (3, 1.0) {$T_0$: 2\,KB};

  %% T1: ops 1-7, offset 2KB, size 1.5KB
  \fill[green!15] (1.1, 2.05) rectangle (7.9, 3.55);
  \draw[green!50!black, thick] (1.1, 2.05) rectangle (7.9, 3.55);
  \node[font=\sffamily\tiny\bfseries, green!55!black] at (4.5, 2.8) {$T_1$: 1.5\,KB};

  %% T2: ops 6-9, offset 0 (reuses T0), size 1KB
  \fill[orange!18] (6.1, 0.05) rectangle (9.9, 0.95);
  \draw[orange!60, thick] (6.1, 0.05) rectangle (9.9, 0.95);
  \node[font=\sffamily\tiny\bfseries, orange!70!black] at (8, 0.5) {$T_2$: 1\,KB};

  %% T3: ops 3-10, offset 3.5KB, size 0.5KB
  \fill[red!12] (3.1, 3.6) rectangle (10.9, 4.1);
  \draw[red!45, thick] (3.1, 3.6) rectangle (10.9, 4.1);
  \node[font=\sffamily\tiny\bfseries, red!60!black] at (7, 3.85) {$T_3$: 0.5\,KB};

  %% Reuse arrow
  \draw[-{Stealth[length=3.5pt]}, dashed, black!40, thick]
    (5.95, 0.5) -- node[above, font=\tiny\itshape, text=black!40] {reuse} (6.05, 0.5);

  %% Pool brace
  \draw[decorate, decoration={brace, amplitude=5pt}, ultra thick, black!25]
    (12.4, 0) -- (12.4, 4.1)
    node[midway, right=7pt, font=\sffamily\scriptsize, text=black!50,
         align=center] {pool\\4.1\,KB};

  %% Birth/death labels
  \draw[blue!50, thick, -{Stealth[length=3pt]}] (0.1, -0.4)
    -- node[below, font=\tiny, text=blue!50] {birth} (0.5, -0.4);
  \draw[blue!50, thick, {Stealth[length=3pt]}-] (5.5, -0.4)
    -- node[below, font=\tiny, text=blue!50] {death} (5.9, -0.4);
\end{tikzpicture}
\caption{Sweep-line memory planning.
  Each tensor is assigned a non-overlapping offset within a single contiguous pool
  based on its dataflow lifetime $[\mathit{birth}, \mathit{death}]$.
  After $T_0$ dies at op~5, $T_2$ reuses the same offset.
  One \texttt{cudaMalloc(pool\_bytes)} replaces all per-operation allocations.}
\label{fig:memplan}
\end{figure}

\textbf{Dynamic shapes.}
New \code{MemoryPlan} computed by the background thread and activated via atomic swap at the iteration boundary.
Foreground never blocked.

\textbf{Activation checkpointing.}
The memory plan provides exact per-tensor cost.
For each forward activation needed in backward: if \code{store\_cost / recompute\_cost > threshold}, recompute.
Per-tensor, adapts to available memory.

\subsection{SMT-Based Kernel Synthesis}

\emph{Design.}
Autotuning (Triton, AutoTVM/Ansor) searches configuration space by benchmarking candidates---hundreds to thousands of trials, finding local optima.
Crucible formulates kernel configuration as a constraint optimization problem.

For a given operation, shape, and target device, \code{CostModel.h} defines the complete problem:

\begin{itemize}
  \item \textbf{Axioms} (\code{HardwareProfile}).
    Measured hardware parameters: \code{num\_sms}, \code{warp\_size}, \code{max\_warps\_per\_sm}, \code{regs\_per\_sm}, \code{max\_regs\_per\_thread} (255), \code{smem\_per\_sm}, \code{tmem\_per\_sm}, four-level bandwidth hierarchy (\code{smem\_bw\_per\_sm}, \code{l2\_bw}, \code{hbm\_bw}), corresponding latencies, peak throughput per dtype (fp64 through fp4), and \code{launch\_ns} (${\sim}3$\,\textmu{}s).
    Presets exist for B200 (sm\_100), H100 (sm\_90), MI300X (gfx942), and A100 (sm\_80).
    Longitude replaces presets with calibrated measurements.
  \item \textbf{Variables} (\code{KernelConfig}).
    \code{tile\_m}, \code{tile\_n}, \code{tile\_k} (\code{uint16}), \code{pipeline\_stages} (1--7), \code{warps\_per\_block} (1--32), \code{smem\_bytes}, \code{regs\_per\_thread}, \code{vec\_width}.
  \item \textbf{Constraints} (C1--C8).
    \code{smem\_bytes} $\leq$ \code{smem\_per\_sm}.
    \code{regs\_per\_thread} $\leq$ \code{max\_regs\_per\_thread}.
    \code{warps\_per\_block} $\geq$ 1.
    \code{warps\_per\_block} $\times$ \code{warp\_size} $\leq$ 1024.
    Positive tile dimensions.
    \code{pipeline\_stages} $\in [1,7]$.
    Each maps mechanically to a Z3 \code{(assert ...)}.
  \item \textbf{Objective.}
    Minimize roofline execution time: $\max(\text{compute\_ns}, \text{memory\_ns}) + \text{launch\_ns}$.
    Compute: $\text{FLOPs} / (\text{peak\_tflops} \times 10^3)$.
    Memory: $\text{bytes} / \text{hbm\_bw}$.
    The ridge point ($\text{peak\_tflops} \times 10^3 / \text{hbm\_bw}$) determines whether a kernel is compute-bound or memory-bound.
\end{itemize}

A Z3 \code{optimize} query finds the configuration minimizing the objective subject to all constraints.
The solution is optimal within the analytical model (Figure~\ref{fig:roofline}).

\begin{figure}[ht]
\centering
\begin{tikzpicture}[xscale=1.1, yscale=0.9]
  %% Axes
  \draw[black!40, -{Stealth[length=4pt]}] (0,0) -- (7.5,0)
    node[right, font=\sffamily\scriptsize]
    {operational intensity (FLOP/byte)};
  \draw[black!40, -{Stealth[length=4pt]}] (0,0) -- (0,5)
    node[above, font=\sffamily\scriptsize, align=center]
    {throughput\\(TFLOPS)};

  %% Ridge point
  \pgfmathsetmacro{\ridge}{3.5}
  \pgfmathsetmacro{\peak}{4.0}

  %% Memory-bound slope
  \draw[blue!50, ultra thick] (0,0) -- (\ridge, \peak);
  %% Compute-bound ceiling
  \draw[red!50, ultra thick] (\ridge, \peak) -- (7.2, \peak);

  %% Ridge point marker
  \fill[black!60] (\ridge, \peak) circle (2.5pt);
  \node[font=\sffamily\tiny, text=black!50, above right=2pt]
    at (\ridge, \peak) {ridge point};

  %% Region labels
  \node[font=\sffamily\scriptsize, text=blue!55, rotate=38] at (1.3, 1.5)
    {memory-bound};
  \node[font=\sffamily\scriptsize, text=red!50] at (5.5, 4.3)
    {compute-bound};

  %% Example kernels
  \fill[orange!70] (1.2, 1.35) circle (3pt)
    node[right=4pt, font=\sffamily\tiny, text=black!55] {LayerNorm};
  \fill[orange!70] (2.0, 2.25) circle (3pt)
    node[right=4pt, font=\sffamily\tiny, text=black!55] {Softmax};
  \fill[green!60!black] (4.5, \peak) circle (3pt)
    node[below right=1pt, font=\sffamily\tiny, text=black!55] {GEMM 4096};
  \fill[green!60!black] (5.8, \peak) circle (3pt)
    node[below right=1pt, font=\sffamily\tiny, text=black!55] {Conv2D};
  \fill[purple!60] (3.2, 3.4) circle (3pt)
    node[left=4pt, font=\sffamily\tiny, text=black!55] {SDPA};

  %% Z3 annotation
  \draw[decorate, decoration={brace, amplitude=4pt}, thick, black!25]
    (7.4, 0) -- (7.4, \peak)
    node[midway, right=7pt, font=\sffamily\tiny, text=black!45,
         align=left, text width=5.5em]
    {Z3 optimizes\\tile, pipeline,\\warps to reach\\roofline ceiling};

  %% HBM BW label
  \node[font=\sffamily\tiny, text=blue!40, rotate=90] at (-0.5, 2)
    {HBM bandwidth};
  %% Peak label
  \node[font=\sffamily\tiny, text=red!40] at (0.8, 4.3)
    {peak TFLOPS};
\end{tikzpicture}
\caption{Roofline model for SMT kernel synthesis.
  Hardware parameters (peak TFLOPS, HBM bandwidth) define the roofline.
  Each kernel's operational intensity determines whether it is
  memory-bound or compute-bound.
  Z3 finds the tile/pipeline/warp configuration that reaches the
  roofline ceiling, subject to hardware constraints C1--C8.
  Longitude calibration adjusts the roofline to measured reality.}
\label{fig:roofline}
\end{figure}

\textbf{Qualification.}
The roofline model does not capture L1/L2 cache behavior, instruction scheduling, warp divergence, or TLB pressure.
Longitude calibration (Section~6) computes per-kernel-class correction factors.
Formal safety guarantees---no out-of-bounds access, no bank conflicts, coalesced access---are model-independent and hold on actual hardware.

\textbf{Scope.}
Applies to any kernel with bounded discrete configuration: GEMM (tile/pipeline/warp), convolution (tile/implicit\_gemm), attention (tile/pipeline/TMA), normalization (vector\_width/rows\_per\_block), reduction (block\_dim/strategy), collective communication (algorithm/channels/chunk\_size).

\subsection{Graduated Divergence Detection}

Four severity levels based on which hash component diverges:

\begin{enumerate}
  \item \textbf{SchemaHash mismatch}---different operation. Hard diverge: immediate eager fallback, re-record.
  \item \textbf{ShapeHash mismatch}---same operation, different shapes. Hard diverge: kernels and memory plan invalid. Fallback, re-record, recompile.
  \item \textbf{ScopeHash mismatch}---same operation, same shapes, different module context (code refactoring). Soft warning: compiled plan remains valid.
  \item \textbf{CallsiteHash mismatch}---same shapes, same module, different source location. Softest warning: logged only.
\end{enumerate}

Figure~\ref{fig:divergence} illustrates the graduated detection hierarchy.

\begin{figure}[ht]
\centering
\begin{tikzpicture}[
  level/.style={draw=black!30, rounded corners=2pt,
    minimum width=17em, minimum height=2em,
    font=\sffamily\small, align=center,
    blur shadow={shadow blur steps=3, shadow xshift=0.3pt,
                 shadow yshift=-0.3pt, shadow blur radius=1pt}},
  arr/.style={-{Stealth[length=4pt]}, thick, black!30},
  chk/.style={font=\sffamily\scriptsize, text=black!50},
]
  \node[level, fill=red!8] (s1)
    {\textbf{Level 1:} SchemaHash $\neq$ \quad {\scriptsize different op}};
  \node[level, fill=orange!8, below=0.6em of s1] (s2)
    {\textbf{Level 2:} ShapeHash $\neq$ \quad {\scriptsize different shapes}};
  \node[level, fill=yellow!10, below=0.6em of s2] (s3)
    {\textbf{Level 3:} ScopeHash $\neq$ \quad {\scriptsize different module}};
  \node[level, fill=green!6, below=0.6em of s3] (s4)
    {\textbf{Level 4:} CallsiteHash $\neq$ \quad {\scriptsize different source loc}};

  %% Severity labels
  \node[chk, left=1em of s1] {\textcolor{red!60}{\textbf{hard}}};
  \node[chk, left=1em of s2] {\textcolor{orange!70}{\textbf{hard}}};
  \node[chk, left=1em of s3] {\textcolor{yellow!70!black}{\textbf{soft}}};
  \node[chk, left=1em of s4] {\textcolor{green!55!black}{\textbf{log}}};

  %% Action labels
  \node[chk, right=0.8em of s1, text width=8em, align=left]
    {fallback to eager\\re-record, recompile};
  \node[chk, right=0.8em of s2, text width=8em, align=left]
    {fallback to eager\\kernels + plan invalid};
  \node[chk, right=0.8em of s3, text width=8em, align=left]
    {warning only\\plan remains valid};
  \node[chk, right=0.8em of s4, text width=8em, align=left]
    {logged only\\no action};

  %% Check flow
  \draw[arr] ([yshift=1.2em]s1.north) -- (s1.north)
    node[above right=-0.1em and 0.2em, chk] {compiled op $i$};
  \draw[arr] (s1.south) -- (s2.north)
    node[midway, right=0.3em, chk] {match $\checkmark$};
  \draw[arr] (s2.south) -- (s3.north)
    node[midway, right=0.3em, chk] {match $\checkmark$};
  \draw[arr] (s3.south) -- (s4.north)
    node[midway, right=0.3em, chk] {match $\checkmark$};
\end{tikzpicture}
\caption{Graduated divergence detection.
  Four hash levels are checked in order of decreasing severity.
  A SchemaHash mismatch (different operation) triggers immediate eager fallback.
  A CallsiteHash mismatch (refactored code, same semantics) is logged without action.}
\label{fig:divergence}
\end{figure}

\textbf{Recovery protocol.}
After divergence, the foreground thread sets \code{BackgroundThread::reset\_requested} (atomic bool).
The background thread detects this at the start of each drain cycle, discards accumulated trace and resets the \code{IterationDetector}---preventing leftover signature ops from producing false iteration boundaries and garbage regions.
Recording skips during alignment; the divergent op itself is not recorded.
If the divergence was temporary, the system may re-encounter the original behavior and switch back via BranchNode arm selection.

\subsection{Shadow Handle Execution Model}

\emph{Design.}
In Tier~2, the Vessel interceptor returns \emph{shadow handles}: full tensor objects (\code{ConductorTensorImpl}) with correct metadata (shape, strides, dtype, device) whose storage points into the pre-planned memory pool.
The GPU executes compiled kernels asynchronously on CUDA streams; Python holds shadows and inspects only metadata.

Synchronization points (the only moments Python blocks): \code{.item()}, \code{.cpu()}, \code{.numpy()}, \code{print()}, data-dependent control flow, unrecognized operations.
Everything else returns immediately.
For 1,000 operations with one \code{loss.item()}: 999 shadow returns (${\sim}2$\,\textmu{}s total) + 1 sync (${\sim}10$\,\textmu{}s).

\subsection{Execution Tiers}

Three tiers of increasing optimization:

\textbf{Tier~1 (implemented).}
Eager replay with divergence detection.
Operations replay in recorded order using Vessel dispatch for execution but the static memory plan for allocation.
Guards check hash consistency at each operation.
Validates the recording and divergence infrastructure without custom kernel compilation.

\textbf{Tier~2 (designed).}
Shadow handle execution.
Foreground returns pre-allocated handles; background pre-schedules kernel launches on CUDA streams.
Python and GPU fully decoupled.
Target: ${\sim}2$\,ns per Vessel dispatch (one pointer advance + one hash guard check).

\textbf{Tier~3 (designed).}
CUDA Graph capture and replay.
Entire compiled iteration captured as a CUDA Graph, replayed in one API call.
Target: ${\sim}50$\,ns per iteration, independent of operation count.
Requires static shapes across iterations.

Each tier is a strict superset of the previous tier's correctness infrastructure.
Divergence detection and fallback are identical across tiers.
